\section{Guided Semantic-Patch Generation}
\label{sec:method}

The input is a migration $M$ (its id and the old/new API line patterns used to mine it) and its training siblings:
pairs $(b_i, a_i)$ of one driver file before and after its hand conversion, all committed before any target was
written. The output is one Coccinelle semantic patch for $M$, or an explicit abstention. No step reads a held-out
target. The design follows guided code generation (GCG)~\cite{almorsi2025gcg}: decompose the task into a tree,
generate the leaves bottom-up with tests, and let a critic check the assembled result. Here the ``tests'' are the
training siblings themselves, and the critic is the C compiler. Figure~\ref{fig:pipeline} shows the four stages.

% Pipeline of the guided framework (included by sections/method.tex)
\begin{figure}[t]
\centering
\begin{tikzpicture}[
  node distance=5mm and 5mm,
  box/.style={draw, rounded corners, align=center, font=\footnotesize, minimum height=12mm, text width=30mm},
  det/.style={box, fill=gray!12},
  llm/.style={box, fill=blue!10},
  arr/.style={-{Stealth[length=2mm]}}]
\node[det] (s1) {\textbf{1.~Variant tree}\\hunks $\to$ leaves\\(script)};
\node[llm, right=of s1] (s2) {\textbf{2.~Leaf rules}\\model + \texttt{spatch}\\feedback, $\le 3$ rounds};
\node[det, right=of s2] (s3) {\textbf{3.~Composition}\\greedy, scored on\\all siblings (script)};
\node[llm, right=of s3] (s4) {\textbf{4.~Compile critic}\\build, repair\\or abstain};
\node[det, below=of s2, text width=62mm, minimum height=7mm] (fix)
  {auto-repair $\cdot$ generality critic $\cdot$ lint (script)};
\node[above=2mm of s1, font=\footnotesize] (in) {training siblings};
\node[above=2mm of s4, font=\footnotesize] (out) {one rule, or abstain};
\draw[arr] (in) -- (s1); \draw[arr] (s1) -- (s2); \draw[arr] (s2) -- (s3); \draw[arr] (s3) -- (s4);
\draw[arr] (s4) -- (out);
\draw[arr] ([xshift=-3mm]s2.south) -- ([xshift=-3mm]s2.south |- fix.north);
\draw[arr] ([xshift=3mm]s2.south |- fix.north) -- ([xshift=3mm]s2.south);
\end{tikzpicture}
\caption{The four stages. Shaded boxes are deterministic scripts; blue boxes call the local model. Only the
training siblings are read; held-out targets are never seen.}
\label{fig:pipeline}
\end{figure}


\paragraph{Units and change hunks.} Both versions of a file are split into statement-level units (cut at
\texttt{;}, \texttt{\{}, \texttt{\}} and comment ends; preprocessor lines are units of their own) and normalised for
whitespace. A hunk is a maximal run of removed and added units. All scoring below counts units: a rule's output
$o_i$ on $b_i$ \emph{hits} an expected unit if the unit change also appears between $b_i$ and $o_i$, and makes an
\emph{extra} edit for every unit change between $b_i$ and $o_i$ that the human did not make between $b_i$ and $a_i$.

\subsection{Stage 1: variant tree (decomposition)}
Every hunk of every training sibling becomes a leaf candidate. Hunks that touch a migration API name (the first
identifier of a match of the old or new pattern) are \emph{anchor} hunks and are grouped by the API names they
contain; the remaining hunks are \emph{companion} hunks (for example, a function's return type that changes because
the API it is registered with changed) and are grouped by their C keywords and by identifiers that recur in at least
two training files. Each group is a leaf. Within a leaf, hunks are abstracted to a \emph{shape}: identifiers become
\texttt{ID1}, \texttt{ID2},~\ldots, numbers \texttt{NUM} and strings \texttt{STR}; the \emph{minimal edit} is the
hunk with its common prefix and suffix tokens removed. A leaf is a \emph{task} if it is an anchor leaf or occurs in
at least two training files; companion leaves seen in one file only are treated as possibly unrelated edits.

\subsection{Stage 2: leaf rules with test feedback}
For each task leaf, the model receives a short SmPL primer (a toy \texttt{kmalloc}+\texttt{memset}~$\to$
\texttt{kzalloc} rule and syntax notes, unrelated to any studied migration) and a task prompt with one example per
distinct minimal-edit shape: the enclosing function or struct header (\emph{inside}), the preceding context, the
minimal edit and its abstract pattern. Anchor leaves are generated first. A companion leaf gets an
\emph{interface}: the framework takes the passing rule of the anchor leaf that shares most training files with it,
names its first rule, and writes the header of a dependent rule that inherits all its metavariables; the model only
writes the body. This is GCG's ``a parent sees only its children's interfaces'', inverted for rules that must run
after their anchor.

Each answer goes through four deterministic steps before \texttt{spatch} runs:
(i) \emph{auto-repair} of mechanical SmPL errors (missing or trailing \texttt{@@}, metavariables used as fields or
callees declared as \texttt{identifier}, undeclared \texttt{ID}/\texttt{NUM}/\texttt{STR} placeholders);
(ii) the \emph{generality critic}: a name that the rule matches literally, that the change keeps (it appears on
context lines or on both a removed and an added line), that is not called as a function, and that occurs in only
one training file is driver-specific and is turned into a metavariable;
(iii) a \emph{lint} for rules without \texttt{-}/\texttt{+} lines, comments used in place of code, and unbalanced
brackets; (iv) application to every training file of the leaf. If the rule fails to parse, misses an expected
unit, or makes an extra edit, the model gets feedback for up to $R=3$ rounds: \texttt{spatch}'s error with a
matching generic repair hint, or, per failing file, the missed units with their abstract pattern and the extra edits.
A rule identical to an earlier one raises the sampling temperature by 0.3 for the next round. The best round is
kept (parses, then most hits, then fewest extra edits, then earliest).

\subsection{Stage 3: composition}
Leaf rules are renamed so that rule names cannot collide (a companion rule carries its anchor rule with it) and
composed greedily: candidates are tried in decreasing order of their own hits minus extra edits, and a candidate is
kept only if the composed patch still parses and its score over \emph{all} training files of the migration
strictly increases. Expected units are those of all task leaves, so a leaf that also helps another leaf is
rewarded and a leaf whose edits clash with another is dropped.

\subsection{Stage 4: compile critic}
The composed rule is applied to the most recent training files it changes (up to three whose human version
compiles), and each output is compiled in a worktree at that training commit's parent (x86\_64
\texttt{allmodconfig}, \texttt{CONFIG\_WERROR} off). If an output does not compile, the model gets the compiler
errors, the converted code around them, and the diff between its output and the human version of that training
file, and writes one additional dependent rule through the same interface mechanism (two rounds). Among the rules
whose checked outputs compile (the composed rule, the repaired rules, and the composed rule minus one leaf), the
one with the best training score is kept; if none compiles, the framework \emph{abstains}. A migration without a
rule is counted as a failure on every target, never silently dropped.
